\documentclass{article}

\textwidth=6.5in
\textheight=8.5in
\voffset=-54pt
\oddsidemargin=0in
\evensidemargin=0in
\setlength{\parskip}{8pt}
\setlength{\parindent}{0pt}

\usepackage{amsmath, amssymb}
\usepackage{ifthen}

\usepackage[inline]{enumitem}
\setlist{topsep=0pt, parsep=8pt}

\usepackage[rgb,table]{xcolor}\convertcolorsUfalse
\usepackage{graphicx}

% change hyperref options
\PassOptionsToPackage{hyphens}{url}
\usepackage[bookmarks,colorlinks,linkcolor=black,citecolor=blue,pdfstartview=FitH,urlcolor=blue]{hyperref}
\hypersetup{pdfpagemode=UseNone}
\usepackage{bookmark}

\usepackage{graphicx}
\usepackage{multicol}
\usepackage{framed}
\usepackage{array}

\usepackage{icomma}

\usepackage{soul}

\usepackage{tikz}
\usetikzlibrary{
  matrix,%
  % enables the snake paths
  decorations.pathmorphing,%
  % enables bigger arrows
  decorations.markings,%
  decorations.pathreplacing,%
  decorations.text,%
  calc,%
  patterns,%
  shapes.geometric,%
  arrows,
  decorations.shapes,
  positioning,
  plotmarks,
  intersections
}
\tikzset{>=latex} % sets default arrow type in tikz
\tikzset{font=\normalsize}
\tikzset{mark size=1.5pt, mark options=thin}
\tikzset{pin distance=4pt,
  every pin edge/.style={<-, thin, shorten <= -2pt}}

\interfootnotelinepenalty=10000
\renewcommand{\thefootnote}{(\arabic{footnote})}

\usepackage{multirow, bigdelim}

% change to \usepackage[sols]{handout} to see solutions
\usepackage[sols]{handout}
\renewcommand{\workshopnote}{CA Notes.}    

\newcommand\iscurrentenumi[1]{%
  \ifthenelse{\equal{#1}{\theenumi}}{}{#1}}
\setenumerate[1]{ref=\#\arabic*}
\setenumerate[2]{ref=\protect\iscurrentenumi{\theenumi}(\alph*)}
\newcommand\enumiiprefix[2]{%
  \ifthenelse{{\equal{#1}{\theenumi}}\AND{\equal{#2}{(\alph{enumii})}}}{}{}%
  \ifthenelse{{\equal{#1}{\theenumi}}\AND{\NOT{\equal{#2}{(\alph{enumii})}}}}{#2}{}%
  \ifthenelse{\equal{#1}{\theenumi}}{}{#1#2}%
}
\setenumerate[3]{ref=\protect\enumiiprefix{\theenumi}{(\alph{enumii})}\roman*}

\definecolor{ideacolor}{rgb}{0.7,0,0}
\newcommand{\ideas}[1]{\color{ideacolor} #1 \color{black}}
\newcommand{\ideacolor}{maroon}

\definecolor{purple}{rgb}{0.5,0,1.0}

\usepackage{fancyhdr}
\lhead{\textcolor{white}{This content is protected and may not be shared, uploaded, or distributed.}}
\rhead{}
\rfoot{\vspace*{\fill}\footnotesize\textcopyright{} \emph{Harvard Math Ma}}
\renewcommand{\headrulewidth}{0pt}
\pagestyle{fancy}
\begin{document}

\htitle{Workshop 6}

\begin{workshop}
  Today's metacognitive strategy discussion is about the \emph{spacing effect}: learning is more effective when study sessions are spaced out. Hundreds of experiments have demonstrated the spacing effect in all sorts of learning, from math to foreign language to performing surgery!

  An important thing for students to know about the spacing effect is that spacing out your study sessions \textbf{feels hard}. When you spread out your study sessions, you forget some things in between sessions, so the sessions never feel easy. But that struggle to remember something you're starting to forget is exactly what strengthens your memory of that material!
\end{workshop}

\begin{enumerate}
\item  

  Let $f(x) = \dfrac{x}{x^4 + 48}$.
  \begin{enumerate}
  \item
    \begin{problem}[\vfill]
      Find all critical points of $f$ on $(-\infty, \infty)$.
    \end{problem}

\begin{solution}
    \begin{solution}
    We find critical points by determining where the derivative is 0 or does not exist. Let's rewrite this function first, and apply the Chain Rule:
    $f(x)=\dfrac{x}{x^4+48}=x(x^4+48)^{-1}$
    \begin{align*}
        f'(x)&=1(x^4+48)^{-1}+x(-1(x^4+48)^{-2}(4x^3)\\
        &=\frac{1}{x^4+48}-\frac{4x^4}{(x^4+48)^2}\\
        &=\frac{1}{x^4+48}\cdot \frac{x^4+48}{x^4+48}-\frac{4x^4}{(x^4+48)^2}\\
        &=\frac{x^4+48-4x^4}{(x^4+48)^2}\\
        &=\frac{-3x^4+48}{(x^4+48)^2}
    \end{align*}
    The denominator is always positive, so there are no inputs where the derivative is undefined. Looking only at the numerator, we set it equal to 0 and solve for $x$:
    \begin{align*}
        0&=-3x^4+48\\
       3x^4&=48\\
        x^4&=16\\
        x&= \pm 2
    \end{align*}
    We have two critical points at $x=\pm 2$
\end{solution}
\end{solution}
  \item

    \begin{workshop}
      One thing to make sure students understand is that it isn't necessary to determine whether the critical points are local minima, local maxima, or neither. Since we only care about finding the global maximum, we can just plug the critical points into $f$ and compare with the end behavior.
    \end{workshop}

    \begin{problem}[\nstretch{0.75}]
      Does $f(x)$ have a local maximum on $(-\infty, \infty)$? If so, where? How about a local minimum?
    \end{problem}
\begin{solution}
    Since the denominator in the derivative is always positive, the sign is entirely controlled by the numerator, $48-3x^4$. For vales of $x$ less than -2 and greater than 2, the derivative is negative. For values of $x$ between -2 and 2, the derivative is positive. This means we have a local min at $x=-2$ (the derivative switches from negative to positive there, so the function changes from decreasing to increasing). We also have a local max at $x=2$ (the derivative switches from positive to negative, so the function changes from increasing to decreasing). 
\end{solution}
  \end{enumerate}

  \begin{problemonly}
    We're still working with $f(x) = \dfrac{x}{x^4 + 48}$.
  \end{problemonly}

  \begin{enumerate}[resume]
  \item
    \begin{problem}[\vfill]
      Does $f(x)$ have a global maximum on $[-1, 10]$? If so, where? How about a global minimum?
    \end{problem}

\begin{solution}
    This is where we need to look at the sign chart for $f'$ and what it means about the behavior for $f(x)$ to determine whether these local extrema could be global extrema. \\
    Our local maxima occurs at $x=2$. Considering what we know about where the function is increasing and decreasing, the only place where a greater output might be achieved is the end behavior as $x\to -\infty$. Computing that limit, $\displaystyle \lim_{x\to -\infty} \frac{x}{x^4+48}$, we note the dominant term in the denominator is $x^4$ so $\frac{x}{x^4+48}$ will behave like $\frac{x}{x^4}=\frac{1}{x^3}$ in the limit. 
    $$\displaystyle \lim{x\to -\infty} \frac{x}{x^4+48}=\lim_{x\to-\infty}\frac{1}{x^3} = 0$$
    Comparing this to $f(2)=\frac{2}{2^4+48}$, we see that $f(2)$ is greater than the value of that limit. We conclude the local max is a global max.\\

    Similar reasoning, which compares $\displaystyle \lim_{x\to\infty}f(x)$ to $f(-2)$ confirms $f(-2)$ is a global minimum.
\end{solution}
  

  \end{enumerate}


\item  Let $\displaystyle f(x) = - \frac{2}{(x - 1)^2} + 3$.

  \begin{enumerate}
  \item
    \begin{problem}[\stretch{1.5}]
      Sketch a graph of $f$, and label all asymptotes. 
    \end{problem}

  \item
    \begin{problem}[\vfill]
      Does $f$ achieve a global minimum on $[0, 2]$? If so, where? How about a global maximum? 
    \end{problem}

  \end{enumerate}
\item  \label{fall2023-workshop6-2}

  Jrue is exploring the function $f(x) = \dfrac{x^3 + x^2 - 12 x}{x^2 - 16}$.
  \begin{enumerate}
  \item
    \begin{problem}[1in]
      Where are the discontinuities of $f$? 
    \end{problem}

\begin{solution}
    When a function is described using a single formula, discontinuities will only happen at inputs where the output is undefined. In this function, the only potentially problematic operation is division and we'll only get an undefined result if the denominator is zero. The denominator is zero at $x=\pm 4$. Those are the locations of our two discontinuities. 
\end{solution}
\item
    \begin{problem}[6in]
      Use limits to classify each discontinuity (as a removable discontinuity, jump discontinuity, or vertical asymptote)? 
    \end{problem}

\begin{solution}
    We classify the type of discontinuity by evaluating the limit of the function there.\\
    
    Let's start with $x=4$\\
    
    $\lim_{x\to4^-} \dfrac{x^3 + x^2 - 12 x}{x^2 - 16}$. As $x\to 4$, the numerator $x^3+x^2-12x \to 32$ and the denominator $x^2-16\to 0$. We need to look at one-sided limits.\\
    
    We begin with $\lim_{x\to4^-} \dfrac{x^3 + x^2 - 12 x}{x^2 - 16}$. As $x\to 4^-$, the numerator $x^3+x^2-12x \to 32$ and the denominator $x^2-16\to 0^-$. Therefore $\lim_{x\to4^-} \dfrac{x^3 + x^2 - 12 x}{x^2 - 16}=-\infty$.\\
    
    When just one of the one-sided limits is $\pm \infty$, this is a {\bf vertical asymptote}.\\
    
    Now we look at what's happening as we approach $x=-4$.\\
    
    We begin with $\lim_{x\to-4} \dfrac{x^3 + x^2 - 12 x}{x^2 - 16}$. As $x\to -4$, the numerator $x^3+x^2-12x \to 0$ and the denominator $x^2-16\to 0$. We need to do closer investigation and there looks like potential opportunity to algebra to simplify this function:
    \begin{align*}
        \dfrac{x^3 + x^2 - 12 x}{x^2 - 16}=\dfrac{x(x^2+x-12x)}{x^2-16}=\frac{x(x-3)(x+4)}{(x-4)(x+4)}=\frac{x(x-3)}{x-4}
    \end{align*}\\
    As $x\to -4$, the numerator $x(x-3) \to 28$. The denominator, $x-4\to -8$ 
    The limit exists $\lim_{x\to-4}\dfrac{x^3 + x^2 - 12 x}{x^2 - 16}=\lim{x\to-4} \dfrac{x(x-3)}{x-4} = \frac{28}{-8}$, though the function is undefined there. When the limit exists and the function is undefined there, this is a {\bf removable discontinuity}. \\
\end{solution}
  \end{enumerate}
  Jrue has found that $f$ has two critical points, and he's classified them:
  \begin{itemize}[nosep]
  \item $f(x)$ has a local maximum at $x = 2$
  \item $f(x)$ has a local minimum at $x = 6$
  \end{itemize}
  Now, Jrue is wondering whether $f$ has a global maximum and global minimum on $(-4, 3]$.
  \begin{enumerate}[resume]

  \item
    \begin{problem}[0.75in]
      Is $f$ continuous on $(-4, 3]$? 
    \end{problem}

\begin{solution}

\end{solution}

  \item

    \begin{workshop}
      $f$ is continuous on $(-4, 3]$, and it has only one critical point in this interval. Since that critical point is a local maximum, $f$ must be increasing on $(-4, 2)$ and decreasing on $(2, 3)$, so $f$ must have a global maximum at $x = 2$. There's no more work to do!

      Make sure students are clear that this only works because $f$ has only one critical point in $(-4, 3]$. 
    \end{workshop}

    \begin{problem}[\vfill]
      Does $f$ have a global maximum on $(-4, 3]$? If so, where? What's the least amount of work you can do to answer this? 
    \end{problem}    


  \end{enumerate}
  


\item  \label{pasachoff-4.3.2} % Pasachoff 4.3.2

  A rectangular lot adjacent to a road is to be enclosed by a fence. The side along the road requires reinforced fencing, which costs \$7 per foot. Fencing that costs \$5 per foot can be used for the other three sides. The city overseeing the project has \$3600 to spend on the fence. 

  \begin{enumerate}
  \item
    \begin{problem}[\vfill]
      Express the area of the lot as a function of the length of the side along the road. 
    \end{problem}

    \begin{solution}
      {\textbf{Define variables:} Let $A$ be the area, $x$ be the
        length of the side along the road, and $y$ be the other dimension of
        the lot.}
      \begin{center}
        \begin{tikzpicture}
          \draw (0, 0) -- (2, 0) -- node[right] {$y$} (2, 1) --
          node[above]{$x$} (0, 1) -- cycle;
          \draw (1, 0.5) node {lot};
          \draw (-1, 0) -- (3, 0) node[right] {road};
        \end{tikzpicture}
      \end{center}

      {\textbf{Formula for thing we're trying to maximize:} $A =
        xy$} 

      \ideas{Next, we want to express $A$ as a function of just one
        variable.  To do this, we must relate $x$ and $y$ somehow, by using
        something else we're told in the problem.}  We're told that the cost
      of the lot can be up to \$3600.  If we want to build a lot that is as
      large as possible, we should use all of the money.  So,
      \begin{align*}
        3600 = {} &  \text{cost of fence} \\
        \intertext{Let's break this cost down, since there are two types of
          fence:} %
        = {} &  \text{cost of reinforced fencing} + \text{cost of regular
          fencing} \\
        = {} &  (\text{cost per foot of reinforced fencing})(\text{length of
          reinforced fencing}) \\
        & + (\text{cost per foot of regular
          fencing})(\text{length of regular fencing}) \\
        = {} &  7 x + 5 (x + 2y) \\
        = {} &  12 x + 10 y
      \end{align*}
      We want to express this in terms of just $x$, so we should solve for $y$ in terms of $x$. We get $y = \dfrac{3600 - 12x}{10} = \dfrac{6}{5} (300 - x)$. Plugging this into our formula for $A$ gives
      $A$ in terms of just $x$:
      \[ A(y) = \frac{5}{6} x (300 - x).\]      
    \end{solution}

  \item
    \begin{problem}[1.5in]
      What's the domain of the function you wrote in (a)? 
    \end{problem}

  \item 
    \begin{problem}[\stretch{0.75}]
      Does the function you wrote in (a) have a global maximum on the domain you found in (b)? If so, where? 
    \end{problem}

  \item
    \begin{problem}[\nstretch{0.1}]
      Interpret your answer to (c) in terms of the situation. 
    \end{problem}
  \end{enumerate}


  


\end{enumerate}
\end{document}
